%%%%%%%%%%%%%%%%%%%%%%%%%%%%%%%%%%%%%%%%%%%%%%%%%%%%%%%%%%%%
\section{Mathematik}%
\label{sec:Mathe}
%%%%%%%%%%%%%%%%%%%%%%%%%%%%%%%%%%%%%%%%%%%%%%%%%%%%%%%%%%%%
%
Grundsätzlich gilt alles, was in \parencites{ams1999a}{ams1999b} steht. 
Um eine einheitliche Typografie von
\index{Skalare}Skalaren, \index{Vektoren}Vektoren, \index{Matrizen}Matrizen,
\index{Zufallsvariablen}Zufallsvariablen \etc zu vereinfachen, ist in der Datei
\printfilepathquoted{preambel/Math.tex} eine Menge Kurzkommandos definiert.
In diese Datei einfach mal reinschauen, welche Kurzkommandos es gibt.

Grundsätzlich empfiehlt es sich, eigene Mathe-Makros zu definieren und durchgehend zu verwenden.
Damit kann man bei Bedarf die Formatierung der Variablen an einer einzigen Stelle anpassen,
ohne den gesamten Code des Arbeit durchgehen zu müssen.

Die Einbettung der Formeln geschieht mit der \printkeyword{equation}"=Umgebung.
Da die Formeln normalerweise keine eigenständigen Absätze darstellen,
sollen zwischen dem Fließtext und der Formeldefinition keine leeren Zeilen stehen,
da \LaTeX{} diese als Absätze behandelt und unnötigerweise zusätzliche vertikale Abstände hinzufügt.

Auch sollte man darauf achten, dass zwischen einer Zahl und der physikalischen Einheit ein umbruchgeschütztes Leerzeichen verwendet wird.
Allerdings sollte statt dem üblichen (dehnbaren) umbruchgeschützen Wortzwischenraum, 
welches mit einer Tilde erzeugt wird, ein schmälerer Abstand verwendet werden.
Dieser wird mit \verb+\,+  erzeugt und beträgt standardmäßig 0,1667\,em. % Dies geht so:
%
%\begin{inlinelatex}
%Der eingefügte Abstand beträgt 0,1667\,em.
%\end{inlinelatex}

Auf zwei besondere Kommandos wird näher eingegangen,
weil dies in \LaTeX{} häufig falsch gemacht wird.
\begin{itemize}
  \item Für die Matrixtransponierte gibt es das Kommando \verb#\Tr#, also
	\verb#$A^{\Tr}$# liefert $A^{\Tr}$	
	\item \justifying Bei \index{Integral}Integralen muss das \enquote{Differential-d} gemäß
	ISO in aufrechter Schrift als Operator gesetzt sein mit einem kleinen Abstand
	zum Integranden. Hierfür gibt es das spezielle Kommando \verb#\diff#.
	\Dhe, die Schreibweise \verb#\int^1_0 x^2 d x = \frac{1}{3}# (\cref{eq:Diff-wrong}) ist falsch,
	während \verb#\int^1_0 x^2 \diff x = \frac{1}{3}# (\cref{eq:Diff-correct}) das Richtige liefert:
\end{itemize}

\begin{equation}
	\int^1_0 x^2 d x = \frac{1}{3} \qquad \text{(falsche Typografie!)}
	\label{eq:Diff-wrong}
\end{equation}

\begin{equation}
	\int^1_0 x^2 \diff x = \frac{1}{3} \qquad \text{(richtige Typografie!)}
	\label{eq:Diff-correct}
\end{equation}

